\begin{problem}[33届物理竞赛决赛]{71}
    八、电子偶素原子（Ps）是由电子 $e^{-}$ 与正电子 $e^{+}$ （电子的反粒子，其质量与电子的相同，电荷与电子的大小相等、符号相反）组成的量子束缚体系，其能级可类比氢原子能级得出。根据玻尔氢原子理论，电子绕质子的圆周运动轨道角动量的取值是量子化的，即为 $\hbar$ 的整数倍。考虑到质子质量是有限的，氢原子量子化条件应修正为：电子与质子质心系中相对其质心的总轨道角动量取值为 $\hbar$ 的整数倍。这一量子化条件可直接推广到其它两体束缚体系，如电子偶素等。以下计算结果均保留四位有效数字。
    \begin{subproblem}
        写出电子偶素基态能量。
    \end{subproblem}
    \begin{subproblem}
        电子偶素基态原子不稳定，可很快发生湮没而生成两个光子：
        \begin{equation*}
            \mathrm{Ps} \to \gamma_{1} + \gamma_{2}
        \end{equation*}
        当基态电子偶素原子相对于实验室参照系以远小于光速的某速度运动时发生湮没，在相对于该速度方向的偏角为 $\theta_{1}$ 的方向上观测到生成的一个光子 $\gamma_{1}$ ，同时在相对于光子 $\gamma_{1}$ 速度反方向的偏角为 $\Delta\theta$ （$\Delta\theta << 1$ ）的方向上观测到另一个光子 $\gamma_{2}$ ，如图所示。
        \begin{subsubproblem}
            求基态电子偶素原子速度 $v_{0}$ 的大小、两个光子 $\gamma_{1}$ 和 $\gamma_{2}$ 的能量 $E_{1}$ 和 $E_{2}$ 的表达式；并给出当 $\theta_{1}=\frac{\pi}{3}$ 、$\Delta\theta=3.464\times10^{-3}$ 时 $v_{0}$ 、$E_{1}$ 和 $E_{2}$ 的值。
        \end{subsubproblem}
        \begin{subsubproblem}
            当基态电子偶素原子静止时，求两个光子 $\gamma_{1}$ 、$\gamma_{2}$ 各自的能量 $E_{1}$ 、$E_{2}$ 与单个自由电子静能之差。
        \end{subsubproblem}
    \end{subproblem}
    \begin{subproblem}
        当基态电子偶素原子相对于实验室参照系以与光速 $c$ 可比拟的速度 $v_{0}$ 运动时湮没生成两个光子，求生成的两个光子 $\gamma_{1}$ 和 $\gamma_{2}$ 的能量 $E_{1}$ 和 $E_{2}$ 、以及光子 $\gamma_{2}$ 的运动方向相对于 $v_{0}$ 的方向的偏角 $\theta_{2}$ （如图所示）与 $\theta_{1}$ 之间的关系式；并给出当 $v_{0} = \frac{c}{2}$ 、$\theta_{1} = \frac{\pi}{3}$ 时 $E_{1}$ 、$E_{2}$ 和 $\theta_{2}$ 的值。

        已知：氢原子基态能量 $E_{n=1}^{H} = -13.60\,\mathrm{eV}$ ，电子质量 $m_{e} = 0.5110\,\mathrm{MeV}/c^{2}$ ，质子与电子的质量之比为1836。
    \end{subproblem}
\end{problem}